% A Normal Form Preserving Minimality: An Order-4 Counterexample to the Komori--Kashima Problem
% Studia Logica house style: for submission, move the content into the
% Springer template used by Studia Logica (sections numbered, results
% numbered within sections as in Nakamura--Matsuda 2021, "Presented by"
% line added by the editor, author addresses at the end).
\documentclass[11pt]{article}
\usepackage{amsmath,amssymb,amsthm}
\usepackage[margin=30mm]{geometry}
\usepackage{xcolor}
\usepackage{longtable}
\newtheorem{theorem}{Theorem}[section]
\newtheorem{lemma}[theorem]{Lemma}
\newtheorem{proposition}[theorem]{Proposition}
\newtheorem{corollary}[theorem]{Corollary}
\theoremstyle{definition}
\newtheorem{definition}[theorem]{Definition}
\newtheorem{example}[theorem]{Example}
\newtheorem{remark}[theorem]{Remark}
\newtheorem{problem}[theorem]{Open Problem}

\newcommand{\plot}[1]{\par\noindent\textcolor{blue}{\textsf{[Plot: #1]}}\par}
\newcommand{\IL}{\mathbf{IL}}
\newcommand{\CL}{\mathbf{CL}}
\newcommand{\LC}{\mathbf{LC}}
\newcommand{\Ord}{\mathrm{Ord}}
\newcommand{\sub}{\leqslant_{\mathrm{sub}}}
\newcommand{\ssub}{<_{\mathrm{sub}}}
\newcommand{\St}{\mathsf{S}}
\newcommand{\repl}[3]{#1[#2\rightsquigarrow #3]}


\title{A Normal Form Preserving Minimality: An Order-4
  Counterexample to the Komori--Kashima Problem}
\author{Yoshiki Nakamura \and Naosuke Matsuda}
\date{}

\begin{document}
\maketitle

\begin{abstract}
A formula is minimal in a logic $L$ if it belongs to $L$ and no strictly
more general formula, in the sense of the substitution preorder, belongs to
$L$. We show that a polarized extension-variable translation $\St$,
namely Jeřábek's translation $\varphi^+$ with its ``$\wedge1$'' wrappers
removed, preserves minimality: for every logic $L$
that contains the axiom $I=p\to p$ and $\St(\alpha)$, if $\alpha$ is minimal in
$L$ then so is $\St(\alpha)$. This applies to all implicational
logics containing $\mathbf{BCI}$, to extensions of $\mathbf{FL_e}$ in the
multiplicative language, and to extensions of $\mathbf{FL_{ew}}$
(in particular of intuitionistic logic) in the full language, and in these
cases $L+\alpha=L+\St(\alpha)$. For implicational $\alpha$, $\St(\alpha)$ has
order at most $4$. Consequently, every implicational logic over
$\mathbf{BCI}$ axiomatizable by formulas minimal in classical logic is
axiomatizable by such formulas of order at most $4$, and $4$ is optimal.
Applied to the counter-example of Nakamura and Matsuda to the
Komori--Kashima problem, this yields an explicit order-$4$ axiom of
Gödel--Dummett logic minimal in classical logic, answering their Open
Problem~2. We also show that the base logic $\mathbf{BCI}$ cannot be
weakened to $\mathbf{CK}$ or $\mathbf{BK}$. The implicational results are
formalized in Lean~4.
\end{abstract}

\noindent\textit{Keywords}: Intermediate logic, Substructural logic, Minimal
formulas, Substitution preorder, Axiomatization, Extension variables.

%======================================================================
\section{Introduction}\label{sec:intro}

A formula $\alpha$ is \kl{minimal} in a \kl{logic} $L$~\cite{Kom05} if
$\alpha\in L$ and no formula \kl{strictly more general} than $\alpha$ with
respect to the \kl{substitution preorder} belongs to $L$. For example, Peirce's law
$((a\to b)\to a)\to a$ is \kl{minimal} in classical logic $\CL$, whereas
$(a\to b)\to(a\to b)$ is not, since it is a proper substitution instance of
the classical tautology $p\to p$.

Since \kl{minimality} depends on the syntactic shape of a formula,
transformations of axioms into equivalent ones need not preserve it, and it
is not clear a priori whether an axiomatization by \kl{minimal} formulas
can be brought into a normal form, for instance one of bounded \kl{order},
without losing \kl{minimality}.

This question arises in the problem of Komori and
Kashima~\cite{Kom05,Kas05}, who asked which \kl{logics} can be axiomatized by
formulas \kl{minimal} in $\CL$, and in particular whether every intermediate
logic axiomatizable in this way is either intuitionistic logic $\IL$ or
$\CL$. Nakamura and Matsuda~\cite{NM21} answered the latter question in the
negative: they exhibited a formula $G'$ of \kl{order} $5$ that is \kl{minimal} in
$\CL$ and axiomatizes Gödel--Dummett logic $\LC$ over $\IL$, and they asked
whether a counterexample of \kl{order} $4$ exists \cite[Open Problem~2]{NM21}.
\plot{Check the formulation of the problem against \cite{Kom05},
\cite{Kas05} and \cite[Sect.~1]{NM21}.}

\subsection{Main results}

We show that a translation $\St$ based on \kl{extension variables}
(Definition~\ref{def:S}) preserves \kl{minimality}. It replaces every compound
subformula occurrence $u$ of $\alpha$ by a fresh variable $\ev{u}$ and records
the meaning of $\ev{u}$ in a premise $\dprem{u}$, a one-way definition whose
direction depends on the \kl{polarity} of $u$:
\[
\St(\alpha)=\dprem{u_1}\to\dots\to \dprem{u_m}\to \ev{\rt} .
\]
The translation is a variant of the polarized translation $\varphi^+$ of
Jeřábek~\cite[Def.~3.6]{Jer16}, obtained essentially by removing the
conjunctions ``$\wedge1$'' with which Jeřábek wraps each definition. Since
each premise $\dprem{u}$ contains exactly two implications, $\St(\alpha)$ has
\kl{order} at most $4$, and it can be computed in linear time.

\begin{theorem}[Main Theorem]\label{thm:main}
Let $L$ be a \kl{logic} containing the identity axiom $\axI=a\to a$. If $\alpha$ is
\kl{minimal} in $L$ and $\St(\alpha)\in L$, then $\St(\alpha)$ is \kl{minimal} in
$L$.
\end{theorem}

The hypothesis $\St(\alpha)\in L$ is satisfied for every $\alpha\in L$, and
moreover $L\lplus \alpha=L\lplus \St(\alpha)$, whenever $L$ contains the linear
implicational logic $\BCI$ (Corollary~\ref{cor:many-logics}).
Consequently, $\St$ provides a minimality-preserving normal form for
axioms: over $\BCI$, every axiomatization by $\CL$-minimal formulas
can be replaced by one of \kl{order} at most $4$
(Corollary~\ref{cor:normal-form}). The bound $4$ is optimal, because
classical tautologies of \kl{order} at most $3$ are already intuitionistic
theorems (Proposition~\ref{prop:order3}). Theorem~\ref{thm:main}, in the
generality stated above, and Corollary~\ref{cor:many-logics} have been
formalized in Lean~4 (Appendix~\ref{sec:lean}).

The hypotheses on the base logic cannot be weakened substantially. The
axiom $\axI$ cannot be dropped: there is a \kl{logic} containing $\BC$ but
not $\axI$ in which the conclusion of Theorem~\ref{thm:main} fails
(Proposition~\ref{prop:needI}). Nor can $\BCI$ be replaced by
$\CK$ or $\BK$ in Corollary~\ref{cor:many-logics}
(Appendix~\ref{app:weak}).

\subsection{Application to the Komori--Kashima problem}

Applied to $G'$, the translation yields an explicit axiom $\St(G')$ of
$\LC$ of order $4$ that is \kl{minimal} in $\CL$ (Corollary~\ref{cor:LC}). This
answers \cite[Open Problem~2]{NM21} affirmatively. Since
Corollary~\ref{cor:normal-form} holds over $\BCK$, it also applies
to Komori's original formulation of the problem over $\BCK$
\cite[footnote~1]{NM21}. Moreover, in the study of which logics are
axiomatizable by $\CL$-minimal formulas \cite[Open Problem~1]{NM21}, one
may restrict attention to formulas of order $4$
(Corollary~\ref{cor:normal-form} and Proposition~\ref{prop:order3}).

\subsection{Substructural logics}

The translation and Theorem~\ref{thm:main} extend to languages with
fusion, lattice connectives and constants. The analogue of
Corollary~\ref{cor:many-logics} holds over $\FLe$ in the
multiplicative language and over $\FLew$, in particular over
intuitionistic logic, in the full language. Over $\FLe$ with
additive connectives, however, neither $\St$ nor Jeřábek's $\varphi^+$
serves our purpose (Section~\ref{sec:fl}).

\subsection{Overview of the proof}

The proof of Theorem~\ref{thm:main} is purely syntactic; it uses only
closure under substitution and modus ponens, together with the axiom $\axI$.
It proceeds in two steps. First, a formula of $L$ is \kl{minimal} in $L$ if and
only if no \kl{one-step generalization} of it belongs to $L$, where a
\kl{one-step generalization} is obtained by replacing some occurrences of a
single subformula by a fresh variable (Proposition~\ref{cor:criterion}).
Second, every \kl{one-step generalization} $\alpha'$ of $\St(\alpha)$ can be
decoded into a \kl{one-step generalization} $\beta$ of $\alpha$: there is a
substitution $\theta$ such that
\[
\theta\alpha'=(\gamma_1\to\gamma_1)\to\dots\to(\gamma_k\to\gamma_k)\to\beta
\]
(Lemma~\ref{lem:key}). Hence, if $\alpha'$ belonged to $L$, then so would
$\beta$, contradicting the \kl{minimality} of $\alpha$. Polarity plays no role
in this argument; it is needed only to ensure that $\St(\alpha)\in L$.
Example~\ref{ex:decode} illustrates the decoding for Peirce's law.

\subsection{Related work}\label{sec:related}

The translation $\St$ is built from well-known ingredients. Extension
variables go back to
Tseitin~\cite{Tse68}, and Statman~\cite{Sta79} used them to reduce
intuitionistic propositional logic to its implicational fragment (see also
Švejdar~\cite{Sve03}). One-way definitions chosen according to \kl{polarity}
appear in the Plaisted--Greenbaum translation~\cite{PG86} and in Mints'
normal form~\cite{Min90}, and Rybakov~\cite{Ryb97} uses \kl{extension variables}
in the study of admissible rules.

Our translation is closest to that of Jeřábek~\cite{Jer16}. His translation $\varphi^+$ introduces a fresh
variable for every occurrence of a non-variable subformula, uses the same
polarized one-way definitions, and wraps each of them as $\dprem{u}\wedge1$. The
translation $\St$ differs from $\varphi^+$ only in that it omits these
wrappers, curries the product of the definitions into an iterated
implication, and treats constants as atoms. With the wrappers,
$\FLe\lplus \varphi=\FLe\lplus \varphi^+$ holds for every formula
$\varphi$ in the full language \cite[Thm.~3.8]{Jer16}; Jeřábek uses this to
show that the substructural hierarchy of Ciabattoni, Galatos and
Terui~\cite{CGT08} collapses to the level $\mathcal N_3$ over
$\FLe$. The main point of the present paper is the observation that,
without the wrappers, the translation preserves \kl{minimality}; to our
knowledge, this has not been observed before. The wrappers destroy
\kl{minimality} (Section~\ref{sec:fl}).
\plot{Check the exact forms used in \cite{Sta79}, \cite{Sve03}, \cite{Min90}.}

\subsection{Organization of the paper}

Section~\ref{sec:prelim} fixes notation and terminology,
Section~\ref{sec:translation} defines the translation $\St$, and
Section~\ref{sec:proof} proves Theorem~\ref{thm:main}.
Section~\ref{sec:normal-form} establishes the basic properties of $\St$
over $\BCI$ and derives the normal form, and Section~\ref{sec:kk} applies
it to the Komori--Kashima problem. Section~\ref{sec:remarks} shows that the
axiom $\axI$ cannot be dropped from Theorem~\ref{thm:main}, and
Section~\ref{sec:fl} extends the results to substructural logics.
Section~\ref{sec:conclusion} concludes with open problems.
Appendix~\ref{app:criterion} proves Proposition~\ref{cor:criterion},
Appendix~\ref{sec:lean} describes the Lean formalization, and
Appendix~\ref{app:weak} shows that $\axB$ and $\axC$ cannot be dropped from
Corollary~\ref{cor:many-logics}.

%======================================================================
\section{Preliminaries}\label{sec:prelim}

\paragraph{Formulas and logics.}
We work in a propositional language $\mathcal L$ that contains $\to$ and
possibly some of the binary connectives $\cdot$ (fusion), $\wedge$,
$\vee$ and some constants (among $\bot,\top,0,1$). An \emph{atom} is a
variable or a constant. $\alpha\to\beta\to\gamma$ stands for
$\alpha\to(\beta\to\gamma)$. As in \cite[Sect.~2]{NM21}, a \emph{logic} is a set
$L$ of $\mathcal L$-formulas closed under modus ponens and substitution;
$L+F$ is the smallest logic containing $L\cup F$. $L$ is \emph{consistent}
if it is not the set of all formulas. Since $L$ is closed under
substitution, $L$ is consistent iff no variable belongs to $L$.

We use the following logics (as sets of theorems of their usual Hilbert
calculi, see \cite{GJKO07}):
\begin{itemize}
\item $\mathbf{BCI}$, $\mathbf{BCK}$, $\IL$ and $\CL$ in the implicational
  language. $\IL=\{K,S\}^*$ and $\CL=\IL+P$ as in \cite{NM21}.
\item $\mathbf{FL_e}$ (full Lambek calculus with exchange) and
  $\mathbf{FL_{ew}}$ (with exchange and weakening) in the language
  $\{\to,\cdot,\wedge,\vee,0,1\}$, possibly with $\bot,\top$. Intuitionistic
  logic in the full language contains $\mathbf{FL_{ew}}$.
\end{itemize}
For implicational formulas, $\alpha\in\CL$ iff $\alpha$ is true under every
Boolean valuation.

A \emph{substitution} $\sigma$ maps variables to formulas; $\sigma\alpha$ is
the result of applying it to $\alpha$. We write $\beta\sub\alpha$ if
$\alpha=\sigma\beta$ for some $\sigma$, and $\beta\ssub\alpha$ if
$\beta\sub\alpha$ but not $\alpha\sub\beta$. A formula $\alpha$ is
\emph{minimal in $L$} if $\alpha\in L$ and there is no $\beta\in L$ with
$\beta\ssub\alpha$. The \emph{order} of an implicational formula is given
by $\Ord(p)=1$ and $\Ord(\alpha\to\beta)=\max\{\Ord(\alpha)+1,\Ord(\beta)\}$.

Every logic $L$ satisfies:
\begin{enumerate}
\item[(C1)] if $\gamma\in L$, then $\sigma\gamma\in L$ for every substitution $\sigma$;
\item[(C2)] if $\gamma_1\to\dots\to\gamma_k\to\beta\in L$ and
  $\gamma_1,\dots,\gamma_k\in L$, then $\beta\in L$.
\end{enumerate}

\paragraph{Occurrences.}
We identify an occurrence in a formula $\alpha$ with its position in the
syntax tree: a finite word over $\{0,1\}$. The empty word $\varepsilon$ is
the root. If the subformula at $u$ is $\gamma\circ\delta$ ($\circ$ a binary
connective), then $\gamma$ is at $u0$ and $\delta$ is at $u1$. We write
$\alpha_u$ for the subformula at $u$. We say that $u$ is \emph{strictly
below} $v$ if $v$ is a proper prefix of $u$. An occurrence $u$ is a
\emph{leaf} if $\alpha_u$ is an atom, and a \emph{compound occurrence}
otherwise. Two occurrences of the same formula are never strictly below one
another.

For a set $Q$ of occurrences in $\alpha$, no one of which is strictly below
another, and a variable $z$, let $\alpha[Q:=z]$ be the formula obtained
from $\alpha$ by replacing the subformula at each $u\in Q$ by $z$.

%======================================================================
\section{The Translation}\label{sec:translation}

\AP Fix a formula $\alpha$. We assign a \intro{polarity}, \intro{positive}
or \intro{negative}, to each occurrence in $\alpha$ as follows. The root
$\rt$ is \kl{positive}. If
$\subf{\alpha}{u}=\subf{\alpha}{u0}\to\subf{\alpha}{u1}$, then $u1$ has
the same \kl{polarity} as $u$ and $u0$ the opposite one. \AP For each
\kl{compound occurrence} $u$ we choose a new variable $\intro*\ev{u}$,
called an \intro{extension variable}, so that distinct occurrences receive
distinct variables and no $\ev{u}$ occurs in $\alpha$. For every
occurrence $u$ let
\[
\intro*\rv{u}=\begin{cases}\subf{\alpha}{u},&\text{if $u$ is a \kl{leaf}},\\
\ev{u},&\text{if $u$ is a \kl{compound occurrence}.}\end{cases}
\]
\AP For every \kl{compound occurrence} $u$ we call
$\intro*\body{u}=\rv{u0}\to\rv{u1}$ the \intro{body} of $u$, and we define
the \intro{defining premise} of $u$ by
\[
\intro*\dprem{u}=\begin{cases}\body{u}\to\ev{u},&\text{if $u$ is \kl{positive}},\\
\ev{u}\to\body{u},&\text{if $u$ is \kl{negative}.}\end{cases}
\]

\begin{definition}\label{def:S}
\AP Let $u_1,\dots,u_m$ be the \kl{compound occurrences} of $\alpha$,
listed in post-order; that is, $u_i$ precedes $u_j$ if and only if $u_i$
is \kl{strictly below} $u_j$, or $u_i=w0u'$ and $u_j=w1u''$ for some
words $w,u',u''$. The \intro{translation} of $\alpha$ is the formula
\[
\intro*\St(\alpha)=\dprem{u_1}\to\dots\to\dprem{u_m}\to\rv{\rt} .
\]
\end{definition}

If $\alpha$ is a variable, then $m=0$ and $\St(\alpha)=\alpha$. \AP Let $\intro*\decsub{\alpha}$ be the substitution that maps each
$\ev{u}$ to $\subf{\alpha}{u}$ and fixes all other variables; then
$\decsub{\alpha}(\rv{v})=\subf{\alpha}{v}$ for every occurrence $v$ of
$\alpha$.


In the following examples, each implication node of the syntax tree of
$\alpha$, that is, each \kl{compound occurrence} $u$, is annotated with
$\ev{u}$ and the \kl{polarity} of $u$; the left and right children of $u$
are $u0$ and $u1$.

\begin{example}\label{ex:peirce}
For Peirce's law $\alpha=((a\to b)\to a)\to a$:
\[
\begin{tikzpicture}[syntree,
  level 1/.style={sibling distance=26mm},
  level 2/.style={sibling distance=16mm},
  level 3/.style={sibling distance=9mm}]
\node[ann={$\ev{\rt}^{+}$}] {$\to$}
  child {node[ann={$\ev{0}^{-}$}] {$\to$}
    child {node[ann={$\ev{00}^{+}$}] {$\to$} child {node {$a$}} child {node {$b$}}}
    child {node {$a$}}}
  child {node {$a$}};
\end{tikzpicture}
\qquad
\begin{aligned}
\dprem{00}&=\body{00}\to\ev{00}=(a\to b)\to\ev{00},\\
\dprem{0}&=\ev{0}\to\body{0}=\ev{0}\to(\ev{00}\to a),\\
\dprem{\rt}&=\body{\rt}\to\ev{\rt}=(\ev{0}\to a)\to\ev{\rt}.
\end{aligned}
\]
Hence
\[
\St(\alpha)=((a\to b)\to\ev{00})\to(\ev{0}\to\ev{00}\to a)\to((\ev{0}\to a)\to\ev{\rt})\to\ev{\rt}.
\]
\end{example}

\begin{example}\label{ex:lc}
For the axiom $\alpha=((a\to b)\to c)\to((b\to a)\to c)\to c$ of
Gödel--Dummett logic:
\[
\begin{tikzpicture}[syntree,
  level 1/.style={sibling distance=38mm},
  level 2/.style={sibling distance=19mm},
  level 3/.style={sibling distance=14mm},
  level 4/.style={sibling distance=9mm}]
\node[ann={$\ev{\rt}^{+}$}] {$\to$}
  child {node[ann={$\ev{0}^{-}$}] {$\to$}
    child {node[ann={$\ev{00}^{+}$}] {$\to$} child {node {$a$}} child {node {$b$}}}
    child {node {$c$}}}
  child {node[ann={$\ev{1}^{+}$}] {$\to$}
    child {node[ann={$\ev{10}^{-}$}] {$\to$}
      child {node[ann={$\ev{100}^{+}$}] {$\to$} child {node {$b$}} child {node {$a$}}}
      child {node {$c$}}}
    child {node {$c$}}};
\end{tikzpicture}
\qquad
\begin{aligned}
\dprem{00}&=(a\to b)\to\ev{00},\\
\dprem{0}&=\ev{0}\to(\ev{00}\to c),\\
\dprem{100}&=(b\to a)\to\ev{100},\\
\dprem{10}&=\ev{10}\to(\ev{100}\to c),\\
\dprem{1}&=(\ev{10}\to c)\to\ev{1},\\
\dprem{\rt}&=(\ev{0}\to\ev{1})\to\ev{\rt}.
\end{aligned}
\]
Listing the \kl{defining premises} in this order (post-order), we obtain
\[
\begin{aligned}
\St(\alpha)={}&((a\to b)\to\ev{00})\to(\ev{0}\to\ev{00}\to c)
\to((b\to a)\to\ev{100})\to(\ev{10}\to\ev{100}\to c)\\
&\to((\ev{10}\to c)\to\ev{1})\to((\ev{0}\to\ev{1})\to\ev{\rt})\to\ev{\rt}.
\end{aligned}
\]
\end{example}

%======================================================================
\section{Application to the Komori--Kashima Problem}\label{sec:kk}

In this section we prove Corollaries~\ref{cor:LC} and~\ref{cor:iff}.
We first observe that \kl{order} $3$ is not enough.

\subsection{Formulas of order at most \texorpdfstring{$3$}{3}}

\begin{proposition}\label{prop:order3}
Every formula of $\CL$ of \kl{order} at most $3$ belongs to $\IL$.
\end{proposition}
\begin{proof}
A formula of \kl{order} at most $2$ has the form $p_1\to\dots\to p_k\to q$ with
variables $p_1,\dots,p_k,q$, and we read it as the rule ``from
$p_1,\dots,p_k$ infer $q$''. A formula of \kl{order} at most $3$ has the form
$F=H_1\to\dots\to H_n\to r$, where the $H_i$ are of \kl{order} at most $2$ and
$r$ is a variable. Let $V$ be the closure of the empty set under the rules
$H_1,\dots,H_n$. The valuation that makes exactly the variables in $V$ true
satisfies every $H_i$, so $r\in V$ because $F\in\CL$. Following the
construction of $V$ step by step, we can derive $r$ from $H_1,\dots,H_n$ by
modus ponens in $\IL$. Hence $F\in\IL$.
\end{proof}

Hence a classical tautology of \kl{order} at most $3$ never axiomatizes a
proper extension of $\IL$.

\subsection{An axiom of Gödel--Dummett logic of order \texorpdfstring{$4$}{4}}


\begin{proof}[Proof of Corollary~\ref{cor:LC}]
The formula~\eqref{eq:SGprime} is obtained directly from
Definition~\ref{def:S}. By \cite[Lemma~4.2]{NM21}, $G'$ is \kl{minimal} in $\CL$ and $\IL\lplus G'=\LC$.
By Corollary~\ref{cor:many-logics} applied to $\IL$ and to $\CL$, we have
$\IL\lplus\St(G')=\IL\lplus G'=\LC$, and $\St(G')$ is \kl{minimal} in
$\CL$. We have $\Ord(\St(G'))\le4$ by Lemma~\ref{lem:order}, and
$\Ord(\St(G'))\ge4$ by Proposition~\ref{prop:order3}, since
$\St(G')\in\CL\setminus\IL$.
\end{proof}
\plot{The Lean file checks that the displayed formula equals
\texttt{statman Gprime 8} and has order $4$. The formula~\eqref{eq:SGprime} is generated by \texttt{paper/gen\_SGprime.py}
(supplementary material).}

\subsection{A minimality-preserving normal form}

Corollary~\ref{cor:iff} is the ``only if'' direction of the
following corollary; the ``if'' direction is trivial.

\begin{corollary}\label{cor:normal-form}
Let $L_0$ be a \kl{logic} containing $\BCI$. If $L=L_0\lplus\Gamma$, where
every $\alpha\in\Gamma$ is \kl{minimal} in $\CL$, then
$L=L_0\lplus\{\St(\alpha)\mid\alpha\in\Gamma\}$, and every $\St(\alpha)$ is
\kl{minimal} in $\CL$ and of \kl{order} at most $4$.
\end{corollary}
\begin{proof}
The first two claims follow from Corollary~\ref{cor:many-logics} applied to
$L_0$ and to $\CL$, and the last from Lemma~\ref{lem:order}.
\end{proof}

By Proposition~\ref{prop:order3} and Corollary~\ref{cor:LC}, the bound $4$
cannot be improved. Over $\IL$, Corollary~\ref{cor:normal-form} can be
sharpened slightly: since the axioms of \kl{order} at most $3$ may be
omitted by Proposition~\ref{prop:order3}, a \kl{logic} axiomatizable over
$\IL$ by $\CL$-minimal formulas is either $\IL$ or axiomatizable over $\IL$
by $\CL$-minimal formulas of \kl{order} exactly $4$.

\subsection{Antecedents of order \texorpdfstring{$3$}{3}}\label{sec:order3}

The formula $\St(G')$ is rather long. We show that a counterexample of
order $4$ must have at least two \kl{antecedents} of order $3$, which
partly explains this. For a formula $F=A_1\to\dots\to A_n\to r$ with $r$ a variable, we call
$A_1,\dots,A_n$ the \intro{antecedents} of $F$.

\begin{proposition}\label{prop:one-premise}
Let $F=A_1\to\dots\to A_n\to r\in\CL$, where $r$ is a variable and
$\Ord(F)\le4$. If at most one $A_i$ has order $3$, then either $F\in\IL$
or $\IL\lplus F=\CL$.
\end{proposition}

Note that \kl{minimality} of $F$ is not assumed.

\begin{proof}
As in the proof of Proposition~\ref{prop:order3}, we call a formula
$p_1\to\dots\to p_k\to q$ with variables $p_1,\dots,p_k,q$ ($k\ge0$) a
\intro{Horn clause} and read it as the rule ``from $p_1,\dots,p_k$ infer
$q$''; these are exactly the formulas of \kl{order} at most $2$. Since
$\Ord(F)\le4$, every $A_i$ has \kl{order} at most $3$. Recall that the \kl{antecedents} of a
formula can be permuted in $\IL$.

Let $\Gamma$ be the set of \kl{antecedents} of \kl{order} at most $2$. For a
set $X$ of variables, let $\mathrm{Cl}(X)$ be the closure of $X$ under the rules in
$\Gamma$. We identify a set of variables with the valuation that makes
exactly those variables true. Then:
\begin{enumerate}
\item[(H1)] $\mathrm{Cl}(X)$ satisfies $\Gamma$;
\item[(H2)] if $q\in \mathrm{Cl}(X)$, then $\Gamma,X\vdash_{\IL}q$, by modus ponens
  along the construction of $\mathrm{Cl}(X)$. In particular, if
  $q\in \mathrm{Cl}(\{p_1,\dots,p_k\})$, then $\Gamma\vdash_{\IL} p_1\to\dots\to p_k\to q$.
\end{enumerate}

\proofcase{Case 1: no \kl{antecedent} has \kl{order} $3$.}
Then $F$ is
$\Gamma\to r$. By (H1) and $F\in\CL$, we have $r\in \mathrm{Cl}(\emptyset)$, and hence
$F\in\IL$ by (H2).

\proofcase{Case 2: exactly one \kl{antecedent}, say
$A=B_1\to\dots\to B_m\to s$, has \kl{order} $3$.}
 Then each $B_j$ has \kl{order} at
most $2$. Assume that $F\notin\IL$, and put $V_0=\mathrm{Cl}(\emptyset)$. We first
note two facts.
\begin{itemize}
\item $r\notin V_0$, since otherwise $\Gamma\vdash_{\IL} r$ by (H2).
\item Some $B_j$ is not derivable from $\Gamma$ in $\IL$. Otherwise
  $\Gamma,A\vdash_{\IL} s$. Moreover, the valuation $\mathrm{Cl}(\{s\})$ satisfies
  $\Gamma$ and $A$, so $r\in \mathrm{Cl}(\{s\})$ because $F\in\CL$. Then
  $\Gamma,s\vdash_{\IL} r$ by (H2), and therefore $F\in\IL$, a contradiction.
\end{itemize}
Fix such a $B_j=p_1\to\dots\to p_k\to q$, and put
$V_1=\mathrm{Cl}(\{p_1,\dots,p_k\})$. Then $V_0\subseteq V_1$, and $q\notin V_1$ by
(H2). The valuation $V_1$ falsifies $B_j$ and hence satisfies $A$; it also
satisfies $\Gamma$ by (H1). Since $F\in\CL$, we obtain $r\in V_1$.

Let $H_3=\{0,\tfrac12,1\}$ be the three-element Gödel algebra, in which
$x\Rightarrow y=1$ if $x\le y$ and $x\Rightarrow y=y$ otherwise. Define a
valuation $v$ by $v(t)=1$ if $t\in V_0$, $v(t)=\tfrac12$ if
$t\in V_1\setminus V_0$, and $v(t)=0$ otherwise. A \kl{Horn clause}
$p_1\to\dots\to p_k\to q$ takes the value $1$ under $v$ if and only if
$\min_i v(p_i)\le v(q)$, where the minimum of the empty set is $1$. Now:
\begin{itemize}
\item every clause in $\Gamma$ takes the value $1$, since $V_0$ and $V_1$
  are closed under the rules in $\Gamma$;
\item $B_j$ takes the value $0$, since $v(p_i)\ge\tfrac12$ for all $i$ and
  $v(q)=0$; hence $A$ takes the value $1$;
\item $v(r)=\tfrac12$, since $r\in V_1\setminus V_0$.
\end{itemize}
Therefore $F$ takes the value $\tfrac12$, and $H_3\not\models F$. By
Jankov's criterion \cite[Theorem~1]{Jan68} (see also
\cite[Sect.~3.2]{NM21}), every implicational formula
$F\in\CL\setminus\IL$ with $H_3\not\models F$ satisfies $\IL\lplus F=\CL$.
\end{proof}
\plot{Check that \cite{Jan68} or \cite[Sect.~3.2]{NM21} covers the
implicational fragment.}

\begin{remark}\label{rem:many}
In $\St(\alpha)$, the formula $\dprem{u}$ has order $3$ if $u$ is
\kl{positive}, since $\dprem{u}=\body{u}\to\ev{u}$, and order $2$ if $u$ is
\kl{negative}, since $\dprem{u}=\ev{u}\to \body{u}$. Hence the
\kl{antecedents} of order $3$ in $\St(\alpha)$ are exactly the formulas
$\dprem{u}$ with $u$ a \kl{positive} \kl{non-leaf occurrence} of $\alpha$. By
Proposition~\ref{prop:one-premise}, if $\alpha\in\CL$ has at most one
\kl{positive} implication occurrence, then $\IL\lplus \alpha=\IL\lplus \St(\alpha)$ is
either $\IL$ or $\CL$. The formula $\St(G')$ in~\eqref{eq:SGprime} has $11$ \kl{antecedents}
of \kl{order} $3$, namely the \kl{antecedents} of the form
$(\beta\to\beta')\to\evfp{\gamma}$, where
$\gamma\in\{a\to b,\ b'\to a',\ (a'\to b')\to c'',\ (b\to a)\to c''',\
G'_7,\dots,G'_2,\ G'\}$,
whereas by Proposition~\ref{prop:one-premise} a $\CL$-minimal axiom of
\kl{order} $4$ of a proper intermediate logic must have at least two
\kl{antecedents} of \kl{order} $3$. We do not know the least
possible number.
\end{remark}

%======================================================================
\section{Proof of the Main Theorem}\label{sec:proof}

We follow the two steps described in the introduction.
Section~\ref{sec:criterion} proves the criterion (step~1).
Section~\ref{sec:decoding} proves the decoding (step~2) and
Theorem~\ref{thm:main}.

\subsection{A criterion for minimality}\label{sec:criterion}

\begin{definition}\label{def:onestep}
Let $\delta$ be a formula, let $Q$ be a nonempty set of occurrences of
$\delta$ in $\alpha$ (that is, $\alpha_u=\delta$ for $u\in Q$), and let $z$
be a variable not occurring in $\alpha$. If $\delta$ is a variable, assume
moreover that some occurrence of $\delta$ in $\alpha$ is not in $Q$. Then
$\alpha[Q:=z]$ is called a \emph{one-step generalization} of $\alpha$.
\end{definition}

\begin{example}
For $\alpha=(a\to b)\to(a\to b)$, both $z\to(a\to b)$ (with $\delta=a\to b$)
and $(z\to b)\to(a\to b)$ (with $\delta=a$) are one-step generalizations of
$\alpha$. $(z\to b)\to(z\to b)$ is not one, because all occurrences of $a$
are replaced.
\end{example}

\begin{lemma}\label{lem:strict}
Every one-step generalization $\beta=\alpha[Q:=z]$ of $\alpha$ satisfies
$\beta\ssub\alpha$.
\end{lemma}
\begin{proof}
$\beta\sub\alpha$, because substituting $\delta$ for $z$ in $\beta$ gives
$\alpha$ ($z$ does not occur in $\alpha$). Suppose, for a contradiction, that
$\tau\alpha=\beta$ for some $\tau$. Call connectives and constants
\emph{symbols}. A substitution never decreases the number of occurrences of
symbols.

If $\delta$ is not a variable, it contains a symbol, so $\beta$ has fewer
symbols than $\alpha$. This is a contradiction.

If $\delta$ is a variable $p$, then $\alpha$ and $\beta$ have the same number
of symbols, so $\tau$ maps each variable of $\alpha$ to a variable. Then
$\tau\alpha$ has the same syntax tree as $\alpha$, with the variable
$\tau(p)$ at every leaf where $\alpha$ has $p$. But at these leaves $\beta$
has $z$ (at $u\in Q$) and $p$ (at the occurrence outside $Q$). Hence
$\tau(p)=z$ and $\tau(p)=p$, a contradiction.
\end{proof}

\begin{lemma}\label{lem:criterion}
Let $L$ be a logic and $\alpha\in L$. If no one-step generalization of
$\alpha$ is in $L$, then $\alpha$ is minimal in $L$.
\end{lemma}
\begin{proof}
Let $\gamma\in L$ and $\sigma\gamma=\alpha$. We show $\alpha\sub\gamma$.

Suppose first that $\sigma$ maps the variables of $\gamma$ injectively to
variables. Then there is a substitution $\tau$ with $\tau(\sigma(q))=q$ for
every variable $q$ of $\gamma$, and $\tau\alpha=\tau\sigma\gamma=\gamma$.

Otherwise, there is a variable $q$ of $\gamma$ such that either (i)
$\sigma(q)$ is not a variable, or (ii) $\sigma(q)=\sigma(q')$ is a variable
for some variable $q'\ne q$ of $\gamma$. Let $z$ be a variable not occurring
in $\alpha$, and let $\sigma'$ agree with $\sigma$ except that
$\sigma'(q)=z$. Every occurrence of $q$ in $\gamma$ yields an occurrence of
$\delta:=\sigma(q)$ in $\alpha=\sigma\gamma$. Let $Q$ be the set of these
occurrences; $Q\ne\emptyset$. Then $\sigma'\gamma=\alpha[Q:=z]$. In case (ii),
the occurrences of $q'$ in $\gamma$ yield occurrences of the variable
$\delta$ that are not in $Q$. So $\sigma'\gamma$ is a one-step generalization
of $\alpha$, and $\sigma'\gamma\in L$ by (C1). This contradicts the
hypothesis.
\end{proof}

\begin{corollary}\label{cor:criterion}
Let $L$ be a logic. A formula $\alpha\in L$ is minimal in $L$ if and only if
no one-step generalization of $\alpha$ is in $L$.
\end{corollary}
\begin{proof}
``If'' is Lemma~\ref{lem:criterion}. ``Only if'' follows from
Lemma~\ref{lem:strict}.
\end{proof}


\subsection{Decoding one-step generalizations}\label{sec:decoding}

For an occurrence $u$ of $\alpha$, the \emph{reading} of $u$ is the
occurrence of $r_u$ inside the body $b_v$ of its parent $v$, or the last
formula $r_\varepsilon$ of $\St(\alpha)$ when $u=\varepsilon$. The next lemma
lists the facts about the shape of $\St(\alpha)$ that we use. All of them are
immediate from the definition.

\begin{lemma}\label{lem:shape}
\begin{enumerate}
\item[(a)] For each compound occurrence $u$, $x_u$ occurs in $\St(\alpha)$
  exactly twice: at the reading of $u$, and as the \emph{head} of $E_u$
  (the variable side of the link).
\item[(b)] The occurrences of atoms of $\alpha$ in $\St(\alpha)$ are exactly
  the readings of the leaves of $\alpha$. So for each atom $c$ they
  correspond bijectively to the leaves $u$ with $\alpha_u=c$.
\item[(c)] Every compound subformula occurrence in $\St(\alpha)$ is of
  exactly one of the following kinds:
  \begin{itemize}
  \item a body $b_u$ (both immediate subformulas are atoms);
  \item a link $E_u$ (one immediate subformula is an atom, the other a body);
  \item a \emph{tail} $T_i=E_{u_i}\to\dots\to E_{u_m}\to r_\varepsilon$ with
    $1\le i\le m$ (its left immediate subformula is a link).
  \end{itemize}
  In particular, no tail is a subformula of a link, and distinct tails are
  distinct formulas.
\item[(d)] If $b_u=b_v$ as formulas, then $\alpha_u=\alpha_v$. If $E_u=E_v$
  as formulas, then $u=v$.
\end{enumerate}
\end{lemma}
\begin{proof}
(a)--(c) follow by inspecting Definition~\ref{def:S}. For the last claims
of (c): the proper subformulas of a link are its body and atoms, and $T_i$
has $m-i+1$ premises. For (d), note that $r_w$ determines $\alpha_w$: if
$r_w$ is an atom of $\alpha$, then $\alpha_w=r_w$; if $r_w=x_{w'}$, then
$w=w'$. So $b_u=b_v$ gives the same main connective,
$\alpha_{u0}=\alpha_{v0}$ and $\alpha_{u1}=\alpha_{v1}$, hence
$\alpha_u=\alpha_v$. A positive link has the form (body)$\to$(variable), and
a negative one (variable)$\to$(body). So $E_u=E_v$ forces equal polarity and
$x_u=x_v$, hence $u=v$.
\end{proof}

We call a one-step generalization $\St(\alpha)[P:=z]$ \emph{normal} if every
position in $P$ is a reading or a body. We first reduce to normal ones.
Throughout, $z$ is a variable not occurring in $\St(\alpha)$.

\begin{lemma}[Normalization]\label{lem:normal}
If some one-step generalization of $\St(\alpha)$ belongs to a logic $L$,
then some normal one-step generalization of $\St(\alpha)$ belongs to $L$.
\end{lemma}
\begin{proof}
Let $\alpha'=\St(\alpha)[P:=z]\in L$, where $P$ is a set of occurrences of
$\delta$. We go through the kinds of $\delta$ given by
Lemma~\ref{lem:shape}. Each step replaces $\alpha'$ by a substitution
instance, which stays in $L$ by (C1).
\begin{itemize}
\item \emph{$\delta$ is an atom of $\alpha$ or a body.} Then every
  occurrence of $\delta$ is a reading, resp.\ a body, by
  Lemma~\ref{lem:shape}(b),(c). So $\alpha'$ is already normal.
\item \emph{$\delta=x_u$.} By Lemma~\ref{lem:shape}(a) and the definition of
  a one-step generalization, $P$ is one of the two occurrences of $x_u$. If
  it is the reading, $\alpha'$ is normal. If it is the head, exchange the
  variables $x_u$ and $z$. The result is $\St(\alpha)$ with the reading of
  $u$ replaced by $z$, which is normal.
\item \emph{$\delta=E_u$.} By Lemma~\ref{lem:shape}(c),(d), $E_u$ occurs
  only once, so $\alpha'$ is $\St(\alpha)$ with the premise $E_u$ replaced
  by $z$. Substitute for $z$ the link $E_u$ with its head replaced by $z$.
  This gives $\St(\alpha)$ with the head of $E_u$ replaced by $z$, which is
  the previous case.
\item \emph{$\delta=T_i$.} By Lemma~\ref{lem:shape}(c), $T_i$ occurs only
  once, so $\alpha'=E_{u_1}\to\dots\to E_{u_{i-1}}\to z$. Substitute for $z$
  the tail $T_i$ with its last formula $r_\varepsilon$ replaced by $z$. This
  gives $\St(\alpha)$ with the reading of $\varepsilon$ replaced by $z$,
  which is normal.\qedhere
\end{itemize}
\end{proof}

\begin{lemma}[Decoding]\label{lem:key}
Let $\alpha'=\St(\alpha)[P:=z]$ be a normal one-step generalization. Let $Q$
be the set of occurrences $u$ of $\alpha$ such that the reading or the body
of $u$ lies in $P$. Then $\beta=\alpha[Q:=z]$ is a one-step generalization of
$\alpha$, and there is a substitution $\theta$ with
\[
\theta\alpha'=(\gamma_1\to\gamma_1)\to\dots\to(\gamma_k\to\gamma_k)\to\beta .
\]
\end{lemma}
\begin{proof}
\emph{$\beta$ is a one-step generalization.} $Q$ is nonempty. All
$\alpha_u$ ($u\in Q$) are the same formula $\delta_0$:
\begin{itemize}
\item if $\delta$ is an atom, then $\delta_0=\delta$;
\item if $\delta=x_w$, then $Q=\{w\}$;
\item if $\delta$ is a body, use Lemma~\ref{lem:shape}(d).
\end{itemize}
In particular, no element of $Q$ is strictly below another. If $\delta_0$ is
a variable, then $\delta=\delta_0$ and $Q$ misses some leaf labelled
$\delta_0$ by Lemma~\ref{lem:shape}(b).

\emph{The substitution.} Call $v$ \emph{hidden} if it is strictly below some
element of $Q$. For $v$ not hidden, write $\beta_v$ for the subformula of
$\beta$ at $v$ ($\beta_v=z$ for $v\in Q$). Let $\theta$ fix $z$ and the
variables of $\alpha$, and put
\[
\theta(x_v)=\begin{cases}
\alpha_v,&\text{if $v$ is hidden, or the reading of $v$ is in $P$},\\
\beta_v,&\text{otherwise.}
\end{cases}
\]
Here is the rule behind this choice: $x_v$ is sent to $\alpha_v$ when only
its head survives, and to $\beta_v$ when its reading survives.

\emph{Readings.} If $v$ is hidden, $\theta$ maps the reading of $v$ to
$\alpha_v$. If $v$ is not hidden, it maps it to $\beta_v$:
\begin{itemize}
\item if the reading is in $P$, it is $z=\beta_v$;
\item otherwise it is $r_v$, and $\theta(r_v)=\beta_v$ by definition (for a
  leaf $v\notin Q$, $\beta_v=\alpha_v=r_v$).
\end{itemize}
(Positions below an element of $Q$ are not in $P$, because $Q$ has no
nested elements.)

\emph{Premises.} Let $v$ be compound, with main connective $\circ$.
\begin{itemize}
\item If $v$ is hidden, $E_v$ is unchanged, and $\theta$ maps its body to
  $\alpha_{v0}\circ\alpha_{v1}=\alpha_v=\theta(x_v)$.
\item If $v$ is not hidden and $v\notin Q$, $E_v$ is unchanged, and $\theta$
  maps its body to $\beta_{v0}\circ\beta_{v1}=\beta_v=\theta(x_v)$.
\item If $v\in Q$ and its body is in $P$, then $E_v$ has become $z\to x_v$ or
  $x_v\to z$, and $\theta(x_v)=\beta_v=z$.
\item If $v\in Q$ and its reading is in $P$, then $E_v$ is unchanged. Its
  children are hidden, so $\theta$ maps its body to $\alpha_v=\theta(x_v)$.
\end{itemize}
In every case the premise is mapped to $\gamma\to\gamma$ (in either
polarity).

\emph{Conclusion.} The last formula is the reading of $\varepsilon$, which is
not hidden. So it is mapped to $\beta_\varepsilon=\beta$.
\end{proof}

\begin{proof}[Proof of Theorem~\ref{thm:main}]
If $\alpha$ is an atom, then $\St(\alpha)=\alpha$ and there is nothing to
prove. Otherwise, suppose some one-step generalization of $\St(\alpha)$ is in
$L$. By Lemma~\ref{lem:normal}, some normal one-step generalization $\alpha'$
is in $L$. By Lemma~\ref{lem:key} and (C1)--(C3), the one-step
generalization $\beta$ of $\alpha$ is in $L$. This contradicts the
minimality of $\alpha$ (Corollary~\ref{cor:criterion}). Hence no one-step
generalization of $\St(\alpha)$ is in $L$. Since $\St(\alpha)\in L$,
$\St(\alpha)$ is minimal by Corollary~\ref{cor:criterion}.
\end{proof}

Consistency of $L$ is not needed. If $\alpha$ is compound and minimal in $L$,
then the variable $z$, a one-step generalization of $\alpha$, is not in $L$.

\begin{example}\label{ex:decode}
We illustrate Lemma~\ref{lem:key} with Example~\ref{ex:peirce}. Replace the
reading of $00$, i.e.\ the second occurrence of $w=x_{00}$, by $z$:
\[
\alpha'=((a\to b)\to w)\to(y\to(z\to a))\to((y\to a)\to x)\to x .
\]
Here $Q=\{00\}$ and $\beta=(z\to a)\to a$. The substitution is
$w\mapsto a\to b$ (only the head of $w$ survives), $y\mapsto z\to a$ and
$x\mapsto\beta$:
\[
\theta\alpha'=((a\to b)\to(a\to b))\to((z\to a)\to(z\to a))\to
(\beta\to\beta)\to\beta .
\]
Since $\beta\notin\CL$ ($z=0$, $a=0$), neither is $\alpha'$.
\end{example}

%======================================================================
\section{Further Remarks}\label{sec:remarks}

This section collects supplementary results. They are not needed for
the main line of the paper.

\subsection{How far can the base logic be weakened?}\label{sec:weak}
Theorem~\ref{thm:main} needs only $I\in L$ and $\St(\alpha)\in L$. In the
implicational case, the last condition is guaranteed by $\mathbf{BCI}$
(Lemma~\ref{lem:sound}). We show that none of $I$, $B$, $C$ can be
dropped. We write $\mathbf{XY}\dots$ for the implicational logic
axiomatized by the axiom schemes $X,Y,\dots$ among
\[
B=(b\to c)\to(a\to b)\to a\to c,\quad
B'=(a\to b)\to(b\to c)\to a\to c,\quad
C=(a\to b\to c)\to b\to a\to c,
\]
$K=a\to b\to a$ and $I=a\to a$, with modus ponens and substitution.

\paragraph{The axiom $I$ cannot be dropped.}
The proof uses $I$ only to discharge the premises $\gamma\to\gamma$ produced
by the decoding (Lemma~\ref{lem:key}). The next proposition shows that this
use is essential, already over $\mathbf{BC}$.

\begin{proposition}\label{prop:needI}
Let $M=(\{0,1,2\},\to,\{0\})$ with
\[
\begin{array}{c|ccc}
\to&0&1&2\\\hline
0&0&1&1\\
1&0&0&0\\
2&0&1&1
\end{array}
\]
and let $L_M$ be the set of formulas that take the value $0$ under every
valuation in $M$. Then $L_M$ is a logic containing $\mathbf{BC}$ but not $I$.
Let $\alpha=(p\to q)\to(p\to q)$. Then:
\begin{enumerate}
\item[(a)] $\alpha$ is minimal in $L_M$;
\item[(b)] $\St(\alpha)=(x_1\to(p\to q))\to((p\to q)\to x_2)\to((x_1\to x_2)\to
  x_0)\to x_0\in L_M$;
\item[(c)] $\St(\alpha)$ is not minimal in $L_M$: its one-step
  generalization
  \[
  \alpha'=(x_1\to z)\to(z\to x_2)\to((x_1\to x_2)\to x_0)\to x_0 ,
  \]
  obtained by replacing both bodies $p\to q$ by $z$, belongs to $L_M$.
\end{enumerate}
So Theorem~\ref{thm:main} fails for logics not containing $I$, even when
$\St(\alpha)\in L$.
\end{proposition}
\begin{proof}
$L_M$ is closed under substitution. It is closed under modus ponens because
$0\to y=0$ forces $y=0$. A finite check shows that $B$ and $C$ take only the
value $0$, while $2\to2=1$, so $I\notin L_M$.

(a) We have $\alpha\in L_M$: the values of implications lie in $\{0,1\}$,
and $0\to0=1\to1=0$. The one-step generalizations of $\alpha$ are
\[
(z\to q)\to(p\to q),\quad (p\to z)\to(p\to q),\quad (p\to q)\to(z\to q),\quad
(p\to q)\to(p\to z),
\]
\[
z\to(p\to q),\quad (p\to q)\to z,\quad z\to z,\quad z .
\]
Each of them takes a value $\ne0$ under some valuation in $M$. For $z\to z$
take $z=2$, and for $z$ take $z=1$. For each of the remaining six, a
refuting valuation is found by a finite check.

(b) and (c) are finite checks as well. In fact $\alpha'$ is the type of the
linear term $\lambda fgh.\,h(\lambda y.\,g(fy))$, and $\St(\alpha)$ is its
instance under $z\mapsto p\to q$. Both are theorems of $\mathbf{BC}$.
\end{proof}

Decoding $\alpha'$ as in Lemma~\ref{lem:key} gives $\beta=z\to z$. The
premise that would have to be discharged is exactly $z\to z$, which is not
available in $L_M$. Validity of $\alpha$ provides $X\to X$ only for formulas
$X$ whose value is that of an implication.
\plot{Checks: \texttt{tools/bc\_alpha\_matrix.py} (matrix $M$, found with z3
and verified by exhaustive evaluation) and \texttt{tools/bc\_check.py}
(condensed detachment: $\St(\alpha),\alpha'\in\mathbf{BC}$).}


\begin{proposition}\label{prop:weak}
\begin{enumerate}
\item[(a)] \emph{(Without $B$.)} $K$ is minimal in $\mathbf{CK}$, but
  $\St(K)=((b\to a)\to x_1)\to((a\to x_1)\to x_0)\to x_0\notin\mathbf{CK}$.
\item[(b)] \emph{(Without $C$.)} $I$ is minimal in $\mathbf{BB'KI}$, but
  $\St(I)=((a\to a)\to x)\to x\notin\mathbf{BB'KI}$. Likewise $K$ is minimal
  in $\mathbf{BK}$ and in $\mathbf{BB'KI}$, but
  $\St(K)\notin\mathbf{BB'KI}\supseteq\mathbf{BK}$.
\end{enumerate}
Hence, over $\mathbf{CK}$, $\mathbf{BK}$ and $\mathbf{BB'KI}$ (and every
logic between $\mathbf{BK}$ and $\mathbf{BB'KI}$), $\St$ does not preserve
minimality, and $L+\alpha\ne L+\St(\alpha)$ for $\alpha=K$ (resp.\ $I$).
\end{proposition}
\begin{proof}
\emph{Minimality.} The one-step generalizations of $K$ are $z\to b\to a$,
$a\to b\to z$, $a\to z$ and $z$; those of $I$ are $z\to a$, $a\to z$ and
$z$. None of them is classically valid, so none lies in any of the logics
considered, all of which are contained in $\CL$.

(a) Let $M_4=(\{0,1,2,3\},\to,\{0\})$ with the table
\[
\begin{array}{c|cccc}
\to&0&1&2&3\\\hline
0&0&1&2&3\\
1&0&0&1&0\\
2&0&0&0&3\\
3&0&1&2&0
\end{array}
\]
A finite check shows that $M_4$ validates $C$ and $K$ (hence also $I$, which
is derivable in $\mathbf{CK}$), but not $B$. The set $\{0\}$ is closed under
modus ponens, since $0\to y=0$ forces $y=0$. So every $\mathbf{CK}$-theorem
takes the value $0$. With $a=b=1$, $x_1=2$ and $x_0=3$, however,
$b\to a=0$, $(b\to a)\to x_1=2$, $a\to x_1=1$ and $(a\to x_1)\to x_0=0$.
Hence $\St(K)$ takes the value $2\to(0\to3)=2\to3=3\ne0$. Intuitively,
$\St(K)$ needs the composition $\lambda y.\,e_1(\lambda z.\,y)$ inside an
argument, which the combinators $\mathsf C,\mathsf K$ cannot form. By
contrast, $\St(I)=((a\to a)\to x)\to x$ is a $\mathbf{CK}$-theorem (it is the
type of $\mathsf C\mathsf I\mathsf I$).

(b) Let $M_3=(\{0,1,2\},\to,\{0\})$ with $x\to y=0$ if $y\le x$ and
$x\to y=1$ otherwise. $M_3$ validates $B,B',K,I$: this is a finite check.
The set $\{0\}$ is closed under modus ponens, since $0\to y=0$ forces $y=0$.
So every theorem of $\mathbf{BB'KI}$ takes the value $0$ in $M_3$. With
$a=0$ and $x=x_0=2$, however, $\St(I)$ takes the value $(0\to2)\to2=1\to2=1$.
$\St(K)$ is refuted similarly. ($M_3$ refutes $C$ and $W$.)
\end{proof}
\plot{Checks: \texttt{tools/ck\_z3.py} (matrix $M_4$, found with z3 and
verified by exhaustive evaluation), \texttt{tools/bki\_matrix.py} (matrix $M_3$),
\texttt{tools/ck\_cd.py} and \texttt{tools/bk\_cd.py} (condensed-detachment
search, consistent with (a) and (b)). Open: ticket entailment
$\mathbf{T}_\to=\mathbf{BB'IW}$, where $M_3$ does not apply because it
refutes $W$; and whether reordering the links helps. For (b), every order
of links fails already for $\St(I)$, which has a single link.}



\subsection{Additive connectives over $\mathbf{FL_e}$}\label{sec:additive}
\begin{remark}[Additive connectives over $\mathbf{FL_e}$]\label{rem:wrapper}
The proof of Theorem~\ref{thm:main} uses only (C1)--(C3), with
$I\in L$. Polarity is used only in
Lemma~\ref{lem:sound}, which is where $\mathbf{FL_e}$ with $\wedge,\vee$
fails. Two facts, both checked with a cut-free prover for $\mathbf{FL_e}$
(\texttt{tools/fle\_check.py}), show that the problem cannot be removed by
simply adding Jeřábek's ``$\wedge1$'':
\begin{enumerate}
\item[(a)] \emph{Without wrappers, $\St$ need not be sound.}
  $\alpha=(q\to q)\wedge(r\to r)$ is minimal in $\mathbf{FL_e}$, but
  $\St(\alpha)=((q\to q)\to x_0)\to((r\to r)\to x_1)\to((x_0\wedge x_1)\to
  x)\to x$ is not provable in $\mathbf{FL_e}$: each branch of the
  $\wedge$-right rule would have to discard the other branch's link.
\item[(b)] \emph{With wrappers, minimality is always lost.} Jeřábek's
  $\varphi^+$ puts $(E_u\wedge 1)$ in place of $E_u$. The link of the root is
  never discarded, so replacing its $1$ by $z$ gives a one-step
  generalization that is still provable (by $\wedge$-left). Hence
  $\varphi^+$ is not minimal in any $L\supseteq\mathbf{FL_e}$ (including
  $\mathbf{FL_{ew}}$, $\IL$, $\CL$) whenever $\varphi$ is compound. For
  example, $(((p\to p)\to x)\wedge z)\to x$ is provable.
\end{enumerate}
Exhaustive tests over all formulas with at most four binary connectives in
$\to,\wedge,\vee$ and three variables give the following. There are $413$
formulas minimal in $\mathbf{FL_e}$, up to renaming. For $164$ of them
$\St(\alpha)$ is provable, and in all $164$ cases it is minimal, as
Theorem~\ref{thm:main} predicts. $\varphi^+$ is provable in all $413$ cases
and minimal in none. Wrapping only the links strictly below an additive
occurrence gives a provable formula in all $413$ cases, but a minimal one in
only $255$. For $\alpha=p\vee(q\to q)$, for instance, the wrapper can be
removed because the $\vee$-right rule selects one branch. So a
minimality-preserving translation over $\mathbf{FL_e}$, if one exists, must
adapt the wrappers to the actual proof. Whether such a translation exists,
and the non-commutative case $\mathbf{FL}$, are left open.
\end{remark}


\subsection{Remarks on the proofs}

\begin{remark}
Corollary~\ref{cor:criterion} replaces \cite[Props.~3.1--3.2]{NM21}, whose
proofs there are sketched, and it does not need simple formulas.
\end{remark}

\begin{remark}
In Lemma~\ref{lem:sound}(i), linear $\lambda$-terms can be given instead of
sequent derivations. Without exchange the premises of $\St(\alpha)$ cannot be
reordered; see Section~\ref{sec:weak}.
\end{remark}

\subsection{Top-level premises of order 3}\label{sec:order3}

For $F=A_1\to\dots\to A_n\to r$ with $r$ a variable, we call
$A_1,\dots,A_n$ the \emph{top-level premises} of $F$. The next proposition
shows that an order-$4$ counter-example needs at least two top-level
premises of order~$3$. This explains, in part, why $\St(G')$ is long.

\begin{proposition}\label{prop:one-premise}
Let $F=A_1\to\dots\to A_n\to r\in\CL$ with $r$ a variable and
$\Ord(F)\le4$. If at most one $A_i$ has order $3$, then $F\in\IL$ or
$\IL+F=\CL$. Minimality of $F$ is not assumed.
\end{proposition}
\begin{proof}
A formula of order $\le2$ has the form $p_1\to\dots\to p_k\to q$ with
variables $p_i,q$ ($k\ge0$). We call it a \emph{Horn clause} and read it as
the rule ``from $p_1,\dots,p_k$ infer $q$''. Since $\Ord(F)\le4$, every
$A_i$ has order $\le3$. Recall that premises can be permuted in $\IL$.

Let $\Gamma$ be the set of top-level premises of order $\le 2$. For a set
$X$ of variables, let $C(X)$ be the closure of $X$ under the rules in
$\Gamma$. We identify a set of variables with the valuation making exactly
those variables true. Then:
\begin{enumerate}
\item[(H1)] $C(X)$ satisfies $\Gamma$;
\item[(H2)] if $q\in C(X)$, then $\Gamma,X\vdash_{\IL}q$ (modus ponens along
  the closure). In particular, $q\in C(\{p_1,\dots,p_k\})$ implies
  $\Gamma\vdash_\IL p_1\to\dots\to p_k\to q$.
\end{enumerate}

\emph{No premise has order $3$.} Then $F$ is $\Gamma\to r$. By (H1) and
$F\in\CL$, we get $r\in C(\emptyset)$, so $F\in\IL$ by (H2).

\emph{Exactly one premise, $A=B_1\to\dots\to B_m\to s$, has order $3$.}
Then each $B_j$ has order $\le2$. Assume $F\notin\IL$. Put
$C_0=C(\emptyset)$.
\begin{itemize}
\item $r\notin C_0$, since otherwise $\Gamma\vdash_\IL r$ by (H2).
\item Some $B_j$ is not derivable from $\Gamma$ in $\IL$. Otherwise
  $\Gamma,A\vdash_\IL s$. Moreover, the valuation $C(\{s\})$ satisfies
  $\Gamma$ and $A$, so $r\in C(\{s\})$ because $F\in\CL$. Then
  $\Gamma,s\vdash_\IL r$ by (H2), and $F\in\IL$.
\end{itemize}
Fix such a $B_j=p_1\to\dots\to p_k\to q$, and put
$C_1=C(\{p_1,\dots,p_k\})$. Then $C_0\subseteq C_1$, and $q\notin C_1$ by
(H2). The valuation $C_1$ falsifies $B_j$, hence satisfies $A$. It also
satisfies $\Gamma$ by (H1). Since $F\in\CL$, we get $r\in C_1$.

Let $H_3=\{0,\tfrac12,1\}$ be the three-element Gödel algebra, with
$x\Rightarrow y=1$ if $x\le y$ and $x\Rightarrow y=y$ otherwise. Define
$v(t)=1$ if $t\in C_0$, $v(t)=\tfrac12$ if $t\in C_1\setminus C_0$, and
$v(t)=0$ otherwise. A Horn clause $p_1\to\dots\to p_k\to q$ has value $1$
under $v$ iff $\min_i v(p_i)\le v(q)$ (the minimum of an empty set is $1$).
\begin{itemize}
\item Every clause in $\Gamma$ has value $1$, since $C_0$ and $C_1$ are
  closed under the rules in $\Gamma$.
\item $B_j$ has value $0$, since $v(p_i)\ge\tfrac12$ for all $i$ and
  $v(q)=0$. Hence $A$ has value $1$.
\item $v(r)=\tfrac12$, since $r\in C_1\setminus C_0$.
\end{itemize}
So $F$ has value $\tfrac12$, and $H_3\not\models F$. By Jankov's criterion
\cite[Theorem~1]{Jan68} (see \cite[Sect.~3.2]{NM21}), an implicational
$F\in\CL\setminus\IL$ with $H_3\not\models F$ satisfies $\IL+F=\CL$.
\end{proof}

\begin{remark}\label{rem:many}
In $\St(\alpha)$, a positive link $b_u\to x_u$ has order $3$ and a negative
link $x_u\to b_u$ has order $2$. So the top-level premises of order $3$ in
$\St(\alpha)$ are exactly the links of the positive implication occurrences
of $\alpha$. Hence, by Proposition~\ref{prop:one-premise}, if $\alpha\in\CL$
has at most one positive implication occurrence, then
$\IL+\alpha=\IL+\St(\alpha)$ is $\IL$ or $\CL$.

$\St(G')$ has $11$ positive links:
$E_1,E_3,E_5,E_7,E_9,E_{11},E_{12},E_{14},E_{15},E_{17},E_{20}$. By
Proposition~\ref{prop:one-premise}, a $\CL$-minimal axiom of order $4$ for a
proper intermediate logic has at least two top-level premises of order $3$.
The least possible number is open.
\end{remark}

%======================================================================
\section{Formalization}\label{sec:lean}

\plot{Half a page.}
File \texttt{lean/StatmanMinimality.lean} (about 1300
lines, implicational language): Lean~4 core only, no Mathlib, no
\texttt{sorry}, no \texttt{native\_decide}. \texttt{\#print axioms} reports
only \texttt{propext}, \texttt{Classical.choice}, \texttt{Quot.sound}.
Theorem~\ref{thm:main} is formalized in exactly the generality stated: a
logic is a predicate \texttt{Th} with \texttt{Logic Th} (closed under
substitution and modus ponens, proving $\gamma\to\gamma$, consistent), and
\texttt{statman\_minimalIn : Logic Th $\to$ ($\forall$ q, Occ q A $\to$ q <
N) $\to$ MinimalIn Th A $\to$ Th (statman A N) $\to$ MinimalIn Th (statman A
N)}. The classical instance \texttt{statman\_minimal} uses
\texttt{logic\_valid} and \texttt{statman\_valid}. Lemma~\ref{lem:sound}(i)
and Corollary~\ref{cor:equiv} are formalized for every implicational logic
containing $\mathbf{BCI}$ (\texttt{BCIExt Th}: closed under substitution and
modus ponens, proving all instances of $B$, $C$, $I$):
\texttt{BCIExt.statman\_imp : Th (A $\to$ statman A N)},
\texttt{BCIExt.statman\_iff : Th A $\leftrightarrow$ Th (statman A N)}, and
\texttt{BCIExt.statman\_minimal} (Corollary~\ref{cor:many-logics} in the
implicational language, without extra hypotheses). Also formalized:
\texttt{statman\_ord} (order $\le4$) and \texttt{statman\_Gprime}
(Corollary~\ref{cor:LC}: the displayed $\St(G')$ equals
\texttt{statman Gprime 8}). Not formalized: the full language
(Lemma~\ref{lem:sound}(ii),(iii)). Correspondence: \texttt{reduction} and
\texttt{strict} = Corollary~\ref{cor:criterion};
\texttt{decode}/\texttt{decode\_node} = Lemma~\ref{lem:key} and Cases 2--5;
the tail branch of \texttt{statman\_minimalIn} = Case~1.

%======================================================================
\section{Concluding Remarks}\label{sec:conclusion}

We have shown that a polarized extension-variable translation
preserves \kl{minimality} with respect
to the \kl{substitution preorder} in every \kl{logic} that contains the axiom $\axI$
and the translated formula. As a consequence, axiomatizations of
implicational logics over $\BCI$ by $\CL$-minimal formulas can
always be taken to consist of formulas of \kl{order} at most $4$, and
\cite[Open Problem~2]{NM21} has an affirmative answer. We conclude with
some questions that remain open.
\begin{enumerate}
\item Which intermediate logics are axiomatizable by $\CL$-minimal formulas
  \cite[Open Problem~1]{NM21}? By Corollary~\ref{cor:normal-form} and
  Proposition~\ref{prop:order3}, it suffices to consider formulas of order $4$.
\item Is there a minimality-preserving translation over $\FLe$
  with additive connectives, or over the non-commutative calculus
  $\FL$ (Section~\ref{sec:additive})?
\item Does the translation, possibly with the formulas $\dprem{u}$ listed in a different order, preserve
  \kl{minimality} over ticket entailment $\Tto$
  (Appendix~\ref{app:weak})?
\item What is the least number of \kl{antecedents} of order $3$ in a
  $\CL$-minimal axiom of order $4$ of a proper intermediate logic? By
  Proposition~\ref{prop:one-premise} it is at least $2$, whereas $\St(G')$
  has $11$ such \kl{antecedents}.
\end{enumerate}

\paragraph{Acknowledgements.} \plot{JSPS KAKENHI etc.}


\appendix
%======================================================================
\section{The formula $\St(G')$}\label{app:SG}
The implication occurrences of $G'$ are numbered $1,\dots,22$ in pre-order,
and $x_i$ is the extension variable of occurrence $i$. $\St(G')$ lists the
links in the order of the table (post-order), followed by $x_1$. The last
column gives the subformula $G'_u$ at occurrence $i$, so that the decoding
substitution of Lemma~\ref{lem:decode} is $x_i\mapsto G'_u$. The table is
generated by \texttt{paper/gen\_SGprime.py}.

{\small
\begin{longtable}{lllp{5.6cm}}
\caption{The links of $\St(G')$.}\label{tab:SG}\\
$i$ & pol. & $E_i$ & $G'_u$\\\hline
\endfirsthead
$i$ & pol. & $E_i$ & $G'_u$\\\hline
\endhead
$E_{2}$ & $-$ & $x_{2}\to a\to a'$ & $a\to a'$\\
$E_{4}$ & $-$ & $x_{4}\to b'\to b$ & $b'\to b$\\
$E_{6}$ & $-$ & $x_{6}\to c''\to c'$ & $c''\to c'$\\
$E_{8}$ & $-$ & $x_{8}\to c'''\to c'$ & $c'''\to c'$\\
$E_{11}$ & $+$ & $(a\to b)\to x_{11}$ & $a\to b$\\
$E_{10}$ & $-$ & $x_{10}\to x_{11}\to c''$ & $(a\to b)\to c''$\\
$E_{14}$ & $+$ & $(b'\to a')\to x_{14}$ & $b'\to a'$\\
$E_{13}$ & $-$ & $x_{13}\to x_{14}\to c'''$ & $(b'\to a')\to c'''$\\
$E_{18}$ & $-$ & $x_{18}\to a'\to b'$ & $a'\to b'$\\
$E_{17}$ & $+$ & $(x_{18}\to c'')\to x_{17}$ & $(a'\to b')\to c''$\\
$E_{21}$ & $-$ & $x_{21}\to b\to a$ & $b\to a$\\
$E_{20}$ & $+$ & $(x_{21}\to c''')\to x_{20}$ & $(b\to a)\to c'''$\\
$E_{22}$ & $-$ & $x_{22}\to c'\to c$ & $c'\to c$\\
$E_{19}$ & $-$ & $x_{19}\to x_{20}\to x_{22}$ & $((b\to a)\to c''')\to c'\to c$\\
$E_{16}$ & $-$ & $x_{16}\to x_{17}\to x_{19}$ & $((a'\to b')\to c'')\to ((b\to a)\to c''')\to c'\to c$\\
$E_{15}$ & $+$ & $(x_{16}\to c)\to x_{15}$ & $(((a'\to b')\to c'')\to ((b\to a)\to c''')\to c'\to c)\to c$\\
$E_{12}$ & $+$ & $(x_{13}\to x_{15})\to x_{12}$ & $((b'\to a')\to c''')\to (((a'\to b')\to c'')\to ((b\to a)\to c''')\to c'\to c)\to c$\\
$E_{9}$ & $+$ & $(x_{10}\to x_{12})\to x_{9}$ & $((a\to b)\to c'')\to ((b'\to a')\to c''')\to (((a'\to b')\to c'')\to ((b\to a)\to c''')\to c'\to c)\to c$\\
$E_{7}$ & $+$ & $(x_{8}\to x_{9})\to x_{7}$ & $(c'''\to c')\to ((a\to b)\to c'')\to ((b'\to a')\to c''')\to (((a'\to b')\to c'')\to ((b\to a)\to c''')\to c'\to c)\to c$\\
$E_{5}$ & $+$ & $(x_{6}\to x_{7})\to x_{5}$ & $(c''\to c')\to (c'''\to c')\to ((a\to b)\to c'')\to ((b'\to a')\to c''')\to (((a'\to b')\to c'')\to ((b\to a)\to c''')\to c'\to c)\to c$\\
$E_{3}$ & $+$ & $(x_{4}\to x_{5})\to x_{3}$ & $(b'\to b)\to (c''\to c')\to (c'''\to c')\to ((a\to b)\to c'')\to ((b'\to a')\to c''')\to (((a'\to b')\to c'')\to ((b\to a)\to c''')\to c'\to c)\to c$\\
$E_{1}$ & $+$ & $(x_{2}\to x_{3})\to x_{1}$ & $(a\to a')\to (b'\to b)\to (c''\to c')\to (c'''\to c')\to ((a\to b)\to c'')\to ((b'\to a')\to c''')\to (((a'\to b')\to c'')\to ((b\to a)\to c''')\to c'\to c)\to c$\\

\end{longtable}
}


\begin{thebibliography}{99}
\bibitem{CGT08} Ciabattoni, A., N. Galatos, and K. Terui, From axioms to
  analytic rules in nonclassical logics, in \emph{LICS 2008}, IEEE, 2008,
  pp.~229--240.
\bibitem{GJKO07} Galatos, N., P. Jipsen, T. Kowalski, and H. Ono,
  \emph{Residuated Lattices: An Algebraic Glimpse at Substructural Logics},
  Elsevier, 2007.
\bibitem{Jer16} Jeřábek, E., A note on the substructural hierarchy,
  \emph{Mathematical Logic Quarterly} 62(1--2):102--110, 2016.
  arXiv:1507.00700.
\bibitem{Min90} Mints, G., Gentzen-type systems and resolution rules.
  Part~I: propositional logic, in \emph{COLOG-88}, LNCS 417, Springer, 1990,
  pp.~198--231. \plot{check}
\bibitem{PG86} Plaisted, D.~A., and S. Greenbaum, A structure-preserving
  clause form translation, \emph{Journal of Symbolic Computation}
  2(3):293--304, 1986.
\bibitem{Ryb97} Rybakov, V.~V., \emph{Admissibility of Logical Inference
  Rules}, Elsevier, 1997.
\bibitem{Sve03} Švejdar, V., On the polynomial-space completeness of
  intuitionistic propositional logic, \emph{Archive for Mathematical Logic}
  42(7):711--716, 2003.
\bibitem{Tse68} Tseitin, G.~S., On the complexity of derivation in
  propositional calculus, in \emph{Studies in Constructive Mathematics and
  Mathematical Logic, Part II}, 1968, pp.~115--125.
\bibitem{Dum59} Dummett, M., A propositional calculus with denumerable
  matrix, \emph{The Journal of Symbolic Logic} 24(2):97--106, 1959.
\bibitem{Jan68} Jankov, V.~A., On the extension of the intuitionist
  propositional calculus to the classical calculus, and the minimal calculus
  to the intuitionist calculus, \emph{Mathematics of the USSR-Izvestiya}
  2(1):205--208, 1968.
\bibitem{Kas05} Kashima, R., Problems on axiomatization of intermediate
  propositional logics, in \emph{the 39th MLG meeting}, 2005, pp.~59--62.
\bibitem{Kom05} Komori, Y., A problem on logics axiomatized with formulas
  minimal in classical logic, and more (in Japanese), MSJ Autumn Meeting,
  2005.
\bibitem{NM21} Nakamura, Y., and N. Matsuda, On implicational intermediate
  logics axiomatizable by formulas minimal in classical logic: a
  counter-example to the Komori--Kashima problem, \emph{Studia Logica}
  109:1413--1422, 2021.
\bibitem{Sta79} Statman, R., Intuitionistic propositional logic is
  polynomial-space complete, \emph{Theoretical Computer Science}
  9:67--72, 1979.
\bibitem{Lean} de Moura, L., and S. Ullrich, The Lean 4 theorem prover
  and programming language, in \emph{CADE-28}, LNCS 12699, 2021,
  pp.~625--635.
\end{thebibliography}


\end{document}
